\documentclass[
  aps,
  prd,
  reprint,
  superscriptaddress,
  nofootinbib,
  longbibliography,
  floatfix,
  amsmath,
  amssymb
]{revtex4-2}

\usepackage[T1]{fontenc}
\usepackage[utf8]{inputenc}
\usepackage{amsmath,amssymb,amsthm}
\usepackage{graphicx}
\usepackage{bm}
\usepackage[hidelinks]{hyperref}
\usepackage{booktabs}
\usepackage{array}
\usepackage{microtype}
\usepackage{enumitem}
\usepackage{xcolor}

%% ---- Theorem Environments ----
\theoremstyle{plain}
\newtheorem{theorem}{Theorem}
\newtheorem{lemma}[theorem]{Lemma}
\newtheorem{corollary}[theorem]{Corollary}
\newtheorem{proposition}[theorem]{Proposition}
\theoremstyle{definition}
\newtheorem{definition}[theorem]{Definition}
\newtheorem{conjecture}[theorem]{Conjecture}
\theoremstyle{remark}
\newtheorem{remark}[theorem]{Remark}

%% ---- Shortcuts ----
\newcommand{\vp}{\varphi}
\newcommand{\LCDM}{\ensuremath{\Lambda}CDM}
\newcommand{\Lam}{\ensuremath{\Lambda}}
\newcommand{\Sthree}{S^3}
\newcommand{\Reff}{R_{\mathrm{eff}}}
\newcommand{\Rcosm}{R_{\mathrm{cosm}}}
\newcommand{\Fn}{F_n}
\newcommand{\half}{\tfrac{1}{2}}

%% ================================================================
\begin{document}

\title{Compact \texorpdfstring{$S^3$}{S3} Spectral Consistency Analysis of Ring-like Ultra-Large-Scale Structure at \texorpdfstring{$z\approx 0.8$}{z approx 0.8}}

\author{Andrei Preece}
\email{andre.preece@web.de}
\affiliation{Independent Researcher, Berlin, Germany}

\author{Boris Batenin}
\email{batenin@web.de}
\affiliation{Independent Researcher, Berlin, Germany}

\date{April 2026}
\keywords{ultra-large-scale structure, compact topology, $S^3$ spatial sections, Laplacian spectrum, projected-shell law, cosmological principle, Giant Arc, Big Ring, Giant Ring, non-decaying correlations, projection effect}

\begin{abstract}
Recent analyses of Mg\,\textsc{ii}-selected structures at $z\approx 0.8$
report three apparently related ring-like ultra-large-scale features in
a common field: the Giant Arc, the Big Ring, and the Giant Ring. This
paper asks whether a compact, simply connected spatial topology
$\Sthree(\Reff)$ can provide a consistent mode-based description of
such morphology in a linearized setting. We first establish the compact
mode-theory results needed for this interpretation: a no-exponential-tail
lemma for globally supported Laplacian eigenmodes on compact connected
manifolds, a bounded oscillatory two-point function on $\Sthree$, and a
projection result showing how constant-phase shells appear as rings or
arcs on narrow redshift slices. We then introduce an exact
projected-chord size law,
$ d_\ell = 2R_{\rm eff}\sin[\pi/(\ell+1)] $, and show that the earlier
Fibonacci scaling arises as its large-$\ell$ asymptotic limit. Under the
empirical mode assignments $\ell=(1,2,6,8)$ for the Giant Arc, Giant
Ring, Big Ring, and the reported 2D-PSA clustering scale, a single
parameter $R_{\rm eff}=474.511\,\mathrm{Mpc}$ reproduces the four
observed scales at the percent level. A final consistency section states explicitly how the framework treats
spatial topology, the invariant light cone, entropy and the arrow of time,
Bell nonlocality, and Noether conservation on a boundaryless compact slice.
The paper does not claim evidence for compact topology or a derived
selection rule for these mode labels. It is offered as a compact-topology
consistency analysis and as a basis for targeted future tests.
\end{abstract}

\maketitle

%% =================================================================
\section{Introduction}
\label{sec:intro}
%% =================================================================

\subsection{The Observational Context}

The \LCDM\ framework assumes large-scale statistical homogeneity above
the Yadav scale $\lambda_{\mathrm{hom}}\approx 370\,\mathrm{Mpc}$~\cite{Yadav2010}:
the two-point correlation function of density fluctuations is expected
to decay with separation and approach zero for $r\gtrsim\lambda_{\mathrm{hom}}$.

Lopez, Clowes, and Williger~\cite{Lopez2022,Lopez2024,Lopez2026} have
identified three uLSSs at $z\approx 0.8$, traced by Mg\,\textsc{ii}
absorbers in SDSS quasar spectra:
\begin{itemize}[noitemsep]
\item \emph{Giant Arc} (GA)~\cite{Lopez2022}: arc-shaped,
  $\sim 1\,\mathrm{Gpc}$ extent, $(4.5\pm0.6)\sigma$;
\item \emph{Big Ring} (BR)~\cite{Lopez2024}: ring-like,
  $\sim 400\,\mathrm{Mpc}$ diameter, $(4.10\pm0.45)\sigma$;
\item \emph{Giant Ring} (GR)~\cite{Lopez2026}: ring-like,
  $\lesssim 1\,\mathrm{Gpc}$ diameter, two distinct versions at
  $4.73\sigma$ and $4.47\sigma$ (optimum ellipse matching), with
  the field showing $3.53\sigma$ clustering at $\sim 320\,\mathrm{Mpc}$
  in a 2D power spectrum analysis.
\end{itemize}
All three structures occupy the same redshift slice
$z=0.722$--$0.862$ in the same field of view, appearing as a nested
system: the GR surrounds the BR.

Lopez \& Clowes~\cite{Lopez2026} also note that the rings may be 2D
\emph{projections of underlying 3D spheroidal structures}, and that
FLAMINGO-10K fields, under the same 2D power spectrum analysis, are
consistent with random expectations at all relevant scales.

\medskip
The complementary roles of the observational and theoretical
contributions are:
\begin{itemize}[noitemsep]
\item Lopez \& Clowes establish: (i) the existence of ring-like uLSSs
  at $>4\sigma$; (ii) a non-random 2D power spectrum signal in the
  observed field; (iii) consistency of FLAMINGO-10K simulation fields
  with randomness under the same analysis.
\item This paper argues: (i) that such structures can be analysed in a
  compact-topology framework without treating them as immediate
  anomalies; (ii) that the specific discrete mode content of $\Sthree$
  can generate non-decaying correlations in the stated approximation;
  and (iii) that simulations using periodic-box ($T^3$) boundary
  conditions do not directly test the same mode basis.
\end{itemize}

\subsection{What This Paper Does Not Claim}

We state explicitly at the outset:

\begin{description}[noitemsep,labelwidth=0.2em,leftmargin=2em]
\item[\textbf{(N1)}] We do \emph{not} claim that Lopez \& Clowes
  observations \emph{prove} compact $\Sthree$ topology.
\item[\textbf{(N2)}] We do \emph{not} claim that the Fibonacci
  hierarchy is the unique possible spectral coarse-graining across
  all cosmological models.
\item[\textbf{(N3)}] We do \emph{not} claim that \LCDM\ is falsified
  in full; only that its particular topological assumption generates
  a tension with the Lopez--Clowes field.
\item[\textbf{(N4)}] We do \emph{not} claim that every ring-like
  structure in the sky must be topological in origin.
\end{description}

\subsection{Scope of the Paper}

\begin{description}[noitemsep,labelwidth=0.2em,leftmargin=2em]
\item[\textbf{(S1)}] The compact-manifold mode analysis developed
  here implies that pure exponential tail behaviour is not the correct
  generic intuition for globally supported correlations.
\item[\textbf{(S2)}] On $\Sthree$, constant-phase shell structure can
  project to ring-like or arc-like morphology on a narrow redshift
  slice.
\item[\textbf{(S3)}] With an explicit projected-chord law and an empirical mode assignment, the
  observed field can be discussed within a compact-topology consistency
  framework, but only at a phenomenological level.
\end{description}

\medskip
\noindent
\textbf{Scope statement.} The purpose of this paper is limited: it does
not infer a unique cosmic topology, but tests whether a compact,
simply-connected mode basis provides a mathematically controlled way to
organise the reported $z\approx 0.8$ morphology.

\textbf{Logical flow of the argument.}
\begin{itemize}[leftmargin=*,noitemsep]
  \item $\Sthree$ geometry $\Rightarrow$ discrete Laplacian spectrum
    (Lemma~\ref{lem:discrete}).
  \item Discrete global modes $\Rightarrow$ non-decaying global
    correlations (Lemma~\ref{lem:notail};
    Theorem~\ref{thm:nodecay}).
  \item Exact projected-chord law (Proposition~\ref{prop:chordlaw}).
  \item Projected shells on $\Sthree$ $\Rightarrow$ rings on the sky
    (Theorem~\ref{thm:projring}).
  \item empirical one-parameter fit of the observed scales (Sec.~\ref{sec:identification}).
\end{itemize}

\subsection{Organisation}

Section~\ref{sec:spectral} establishes the spectral-geometric
foundations: minimal topology, discrete Laplacian spectrum, and the
No-Exponential-Tail Lemma. Section~\ref{sec:correlations} proves the
main correlation theorem on $\Sthree$. Section~\ref{sec:projection}
derives the Projection Ring Theorem. Section~\ref{sec:fibonacci}
constructs the Fibonacci preferred radii.
Section~\ref{sec:identification} gives the phenomenological
identification with the observed structures. Section~\ref{sec:comparison}
presents the boundary-condition diagnosis of FLAMINGO and lists
testable predictions. Section~\ref{sec:discussion} closes with
open problems and comparison to alternative frameworks.

\medskip\noindent
\textbf{Conventions.} $c=1$ throughout. $\vp=(1+\sqrt5)/2$.
Fibonacci: $F_1=F_2=1$, $F_{n+2}=F_{n+1}+F_n$. Proper distances at
$z=0.802$ for $H_0=67\,\mathrm{km/s/Mpc}$, $\Omega_m=0.315$.

%% =================================================================
\section{Spectral Theory on Compact $S^3$}
\label{sec:spectral}
%% =================================================================

\subsection{Minimal Topological Conditions}

We consider spatial sections $\Sigma$ satisfying:
\begin{enumerate}[label=\textbf{(M\arabic*)},noitemsep]
\item $\Sigma$ is compact, connected, orientable, and boundaryless.
\item $\pi_1(\Sigma)=0$ (trivial fundamental group).
\item No additional topological structure is introduced beyond a round
  Riemannian metric.
\end{enumerate}
By Perelman's proof of the Poincar\'e
conjecture~\cite{Perelman2002,Perelman2003a,Perelman2003b}, the unique
closed 3-manifold satisfying (M1)--(M3) is $\Sigma\cong\Sthree$.

\begin{remark}[Status of $\Reff$ vs global curvature]
  Throughout, $\Reff$ denotes the \emph{effective spectral curvature
  scale} inferred from fitting the ring structures in a narrow redshift
  shell. This is distinct from, and need not equal, the global
  cosmological curvature radius $\Rcosm^{\mathrm{global}}$ constrained
  by Planck 2015~\cite{Planck2016} to exceed $10\,\mathrm{Gpc}$.
  We discuss this two-scale structure in Section~\ref{sec:discussion}.
\end{remark}

\subsection{Geometry of the Three-Sphere}

The round three-sphere $\Sthree(\Reff)$ is the set of points in
$\mathbb{R}^4$ lying at Euclidean distance $\Reff$ from the origin,
equipped with the induced metric. In hyperspherical coordinates
$(\chi,\theta,\phi)$, with $\chi\in[0,\pi]$, the line element is
\begin{equation}
  ds^2 = \Reff^2\left[d\chi^2+
  \sin^2\chi\left(d\theta^2+\sin^2\theta\,d\phi^2\right)\right].
  \label{eq:S3metric}
\end{equation}
The geodesic distance between two points $x,y\in \Sthree$ is
$r=\Reff\,\Theta$, where $\Theta\in[0,\pi]$ is their angular separation.
The total volume is $2\pi^2\Reff^3$.

The scalar Laplace--Beltrami operator $\triangle_0$ is elliptic and
self-adjoint on $L^2(\Sthree,dV)$. Since $\Sthree$ is compact and
connected, its eigenvalues are discrete and its eigenfunctions are
smooth and real-analytic. This is the spectral fact that distinguishes
compact $\Sthree$ from the flat $\mathbb{R}^3$ intuition used in many
large-volume approximations.

\subsection{Discrete Laplacian Spectrum and Global Support}

\begin{lemma}[Discrete global spectrum on $S^3$]
  \label{lem:discrete}
  The scalar Laplacian $\triangle_0$ on $\Sthree(\Reff)$ has purely
  discrete eigenvalues
  \begin{equation}
    \lambda_\ell = \frac{\ell(\ell+2)}{\Reff^2},
    \quad \ell=0,1,2,\ldots,
    \label{eq:eigenvalues}
  \end{equation}
  with degeneracy $g_\ell=(\ell+1)^2$. The corresponding
  hyperspherical harmonics form a complete orthonormal basis of
  $L^2(\Sthree)$.
\end{lemma}

\begin{proof}
  This is the standard scalar spectrum on the round $n$-sphere,
  specialized to $n=3$; see, e.g., Refs.~\cite{Camporesi1994,Berger1971}.
  The degeneracy follows from the dimension of the degree-$\ell$
  harmonic subspace on $S^3$. Equivalently, using
  $\Sthree\simeq SU(2)$, the isometry group is
  $SO(4)\simeq (SU(2)_L\times SU(2)_R)/\mathbb{Z}_2$, and the scalar
  harmonics at level $\ell$ transform in the $(\ell/2,\ell/2)$
  representation of $SU(2)_L\times SU(2)_R$, of dimension $(\ell+1)^2$.
\end{proof}

\begin{remark}[Zero mode and global support]
  For $\ell=0$ one obtains the constant zero mode. For $\ell\geq1$ the
  eigenfunctions are non-constant but globally defined across the whole
  manifold. A nonzero eigenfunction cannot vanish on an open set by the
  unique-continuation principle for elliptic operators; in this precise
  sense the compact modes have global support rather than finite spatial
  support.
\end{remark}

\subsection{Low Multiplets and Four-Dimensional Harmonic Polynomials}
\label{subsec:lowell}

A useful algebraic construction of the scalar harmonics on $S^3$ is obtained
by restricting homogeneous harmonic polynomials on $\mathbb{R}^4$ to the
sphere. Let $u=x/\Reff\in S^3_1$ be the corresponding unit-sphere coordinate.
If $P_\ell(u)$ is homogeneous of degree $\ell$ and harmonic in $\mathbb{R}^4$,
then its restriction to $S^3_1$ is an eigenfunction of the unit-sphere
Laplacian with eigenvalue $\ell(\ell+2)$. Restoring the radius gives
Eq.~\eqref{eq:eigenvalues}.

For orientation, the first multiplets may be written in a real basis as
follows. With the unit-sphere volume convention used in Eq.~\eqref{eq:addition},
\begin{align}
  \ell=0: \quad &Y_0 = (2\pi^2)^{-1/2}, \\
  \ell=1: \quad &Y_{1\mu}(u)=\sqrt{\frac{2}{\pi^2}}\,u_\mu,
    \qquad \mu=1,\ldots,4 .
\end{align}
The $\ell=1$ sector is therefore the four-dimensional dipole multiplet.
The $\ell=2$ sector is represented by traceless quadratic forms
\begin{equation}
  Q_{ij}(u)=u_i u_j-\frac{\delta_{ij}}{4},
  \qquad i,j=1,\ldots,4,
\end{equation}
with nine independent components after imposing tracelessness. Higher
multiplets are obtained analogously from higher-degree harmonic polynomials,
or equivalently from the representation theory of
$SU(2)_L\times SU(2)_R$.

This construction also clarifies the phrase ``global support''. The modes
are global functions on the whole compact manifold; they are not localized
or compactly supported. They may, and generally do, possess nodal sets, so
``global support'' should not be read as ``nowhere zero.''


\subsection{Hyperspherical Harmonics and the Addition Theorem}

Let $\{Y_{\ell\mu}\}_{\mu=1}^{(\ell+1)^2}$ be an orthonormal basis of the
$\ell$th eigenspace. The rotationally invariant projector onto this
eigenspace is given by the addition theorem
\begin{equation}
  \sum_{\mu=1}^{(\ell+1)^2}Y_{\ell\mu}(x)\overline{Y_{\ell\mu}(y)}
  =\frac{\ell+1}{2\pi^2}\,
  \frac{\sin[(\ell+1)\Theta]}{\sin\Theta},
  \label{eq:addition}
\end{equation}
where $\Theta$ is the angular geodesic separation between $x$ and $y$.
We denote the corresponding zonal function by
\begin{equation}
  \Xi_\ell(\Theta):=\frac{\ell+1}{2\pi^2}\,
  \frac{\sin[(\ell+1)\Theta]}{\sin\Theta}.
  \label{eq:zonal}
\end{equation}
This identity is the basic input for the isotropic two-point statistics
used below. The apparent singularity at $\Theta=0$ is removable, with
$\Xi_\ell(0)=(\ell+1)^2/(2\pi^2)$.

\begin{figure}[t]
\centering
\includegraphics[width=\columnwidth]{xi_kernels.png}
\caption{Zonal kernels $\Xi_\ell(\Theta)$ for the low modes used later in
the phenomenological identification. The $\ell=1$ kernel is dipolar, while
higher modes display the increasing sequence of nodes that motivates the
projected-chord scale law.}
\label{fig:Xi}
\end{figure}

\subsection{The Two-Point Correlation Function}

Any square-integrable, mean-zero, statistically isotropic random field
$f$ on $\Sthree(\Reff)$ admits the spectral expansion
\begin{equation}
  f(x)=\sum_{\ell=0}^{\infty}\sum_{\mu=1}^{(\ell+1)^2}
  a_{\ell\mu}Y_{\ell\mu}(x),
  \label{eq:fieldexpansion}
\end{equation}
with uncorrelated coefficients satisfying
$\langle a_{\ell\mu}\overline{a_{\ell'\mu'}}\rangle=
P_\ell\delta_{\ell\ell'}\delta_{\mu\mu'}$. The two-point function depends
only on the angular separation:
\begin{equation}
  C(\Theta):=\langle f(x)f(y)\rangle
  =\sum_{\ell=0}^{\infty}P_\ell\,\Xi_\ell(\Theta).
  \label{eq:corrS3formal}
\end{equation}
At $\Theta=0$, this gives
$C(0)=\sum_\ell P_\ell(\ell+1)^2/(2\pi^2)$, the total variance in the
adopted normalization.

\subsection{The No-Exponential-Tail Lemma}

We now establish the precise formal obstacle to transferring
$\mathbb{R}^3$ intuition to the compact case.

\begin{lemma}[No global pure exponential tail on compact $S^3$]
  \label{lem:notail}
  Let $C:[0,\pi\Reff]\to\mathbb{R}$ be the two-point correlation
  function of a stationary, mean-zero, isotropic random field on
  $\Sthree(\Reff)$ with the spectral form~\eqref{eq:corrS3formal}. If
  the field has nonzero power in any $\ell\geq1$ mode, then the compact
  correlation function cannot be identified with a pure flat-space
  exponential tail
  \begin{equation}
    C(r)=A\,e^{-r/\xi_0},\qquad A>0,\quad \xi_0>0,
    \label{eq:expdecay}
  \end{equation}
  as a globally valid smooth isotropic covariance on the full geodesic
  interval $r\in[0,\pi\Reff]$.
\end{lemma}

\begin{proof}
  A smooth isotropic function on $\Sthree$ is a smooth function of
  $\cos\Theta$, not an arbitrary smooth function of the coordinate
  distance $\Theta$. In particular, regularity at the antipode requires
  the first derivative with respect to $\Theta$ to vanish at
  $\Theta=\pi$ for a smooth zonal covariance. This condition is satisfied
  termwise by the zonal harmonics in Eq.~\eqref{eq:zonal}. A pure
  exponential $A\exp[-\Reff\Theta/\xi_0]$ has derivative
  $-(A\Reff/\xi_0)\exp[-\Reff\Theta/\xi_0]$, which is nonzero at
  $\Theta=\pi$. Hence it cannot be a globally smooth isotropic covariance
  on $S^3$.

  Equivalently, Eq.~\eqref{eq:corrS3formal} expresses any isotropic
  compact covariance as a discrete Gegenbauer/zonal-harmonic expansion.
  Such a compact expansion may approximate exponential decay over a
  finite subinterval, but it cannot equal a flat-space exponential tail
  globally on the compact manifold once nontrivial compact modes are
  present.
\end{proof}

\begin{remark}
  Lemma~\ref{lem:notail} does not state that compact correlations are
  everywhere large, nor does it exclude approximate exponential behaviour
  over a limited observational range. It states only the mathematically
  relevant point for the present paper: the pure exponential tail used as
  flat-space intuition is not a globally valid covariance model on
  compact $\Sthree$ with nontrivial discrete mode content.
\end{remark}

%% =================================================================
\section{Non-Decaying Correlations on $S^3$}
\label{sec:correlations}
%% =================================================================

\begin{theorem}[Compact-mode correlations on $S^3$]
  \label{thm:nodecay}
  Under the assumptions of Lemma~\ref{lem:notail}, the correlation
  function $C(\Theta)$ on $\Sthree(\Reff)$ is a bounded discrete sum (or
  convergent series) of compact zonal harmonics. It is therefore not a
  globally valid pure exponential tail. In particular, any isolated
  nonzero low-$\ell$ contribution remains globally supported on the full
  geodesic interval, including the antipodal regime.
\end{theorem}

\begin{proof}
  Equation~\eqref{eq:corrS3formal} gives $C$ as a discrete expansion in
  the zonal kernels $\Xi_\ell$. For example, the $\ell=1$ contribution is
  \begin{equation}
    P_1\,\Xi_1(\Theta)
    = P_1\cdot\frac{2}{2\pi^2}\cdot\frac{\sin 2\Theta}{\sin\Theta}
    = \frac{2P_1}{\pi^2}\cos\Theta,
    \label{eq:ell1}
  \end{equation}
  which is defined on the full interval $[0,\pi]$ and takes a nonzero
  value at the antipode. Higher-$\ell$ terms behave similarly as global
  compact harmonics with finitely many nodal sets. The assertion then
  follows directly from the spectral representation and from
  Lemma~\ref{lem:notail}.
\end{proof}

\begin{remark}[Physical content]
  Theorem~\ref{thm:nodecay} states that compact $\Sthree$ geometry allows
  globally supported correlation patterns. Coherent structures with
  characteristic geodesic scale $r\sim\pi\Reff/\ell$ may therefore be
  organised as compact eigenmode features, rather than only as accidental
  finite-correlation-length fluctuations.
\end{remark}

\subsection{Contrast with the $\mathbb{R}^3$ Intuition}

On $\mathbb{R}^3$, the Laplacian has continuous spectrum $[0,\infty)$,
no globally-normalisable zero mode, and wave packets decay. Correlation
functions typically satisfy $C(r)\sim e^{-r/\xi_0}$ at large $r$ for
some correlation length $\xi_0$. Within \LCDM\ at late times,
$\xi_0\approx 5$--$10\,\mathrm{Mpc}$ from the matter power spectrum,
and coherent structures of size $\sim 1\,\mathrm{Gpc}$ require
exponentially suppressed probabilities under a finite-correlation-length
assumption. Lemma~\ref{lem:notail} shows that this \emph{decay pattern
is not an option} on compact $\Sthree$ carrying non-trivial mode content.

%% =================================================================
\section{The Projection Ring Theorem}
\label{sec:projection}
%% =================================================================

The observations~\cite{Lopez2026} suggest that the observed rings may be
2D projections of underlying 3D spheroidal structures. We show that
globally-supported spherical shells on $\Sthree$ generically project
to rings or arcs on a narrow redshift slice, providing a structural
link between the geometric framework and the observational morphology.

\begin{theorem}[Projection Ring Theorem]
  \label{thm:projring}
  Let $\mathcal{S}_\rho\subset\Sthree(\Reff)$ be a geodesic sphere of
  radius $\rho$ around a centre $\mathbf{c}$, representing a constant-phase
  shell of a hyperspherical harmonic $Y_\ell^{mm'}$ with
  $\rho\approx \pi\Reff/(\ell+1)$. Let $\pi_{\mathrm{obs}}: \Sthree\to\Sigma_{\mathrm{obs}}$
  be the projection onto a narrow observational redshift slice
  $\Sigma_{\mathrm{obs}}$, transverse to the line of sight of thickness
  $\Delta r$. Then:
  \begin{enumerate}[label=\textup{(\roman*)}]
  \item If $\mathbf{c}\notin\Sigma_{\mathrm{obs}}$, the projection
    $\pi_{\mathrm{obs}}(\mathcal{S}_\rho)$ is a circle or ellipse
    on $\Sigma_{\mathrm{obs}}$ of projected radius
    $r_{\mathrm{proj}} = \sqrt{\rho^2 - h^2}$, where $h$ is the geodesic
    distance from $\mathbf{c}$ to $\Sigma_{\mathrm{obs}}$.
  \item If $|h-\rho|<\Delta r$, the projection degenerates to an
    arc (partial ring), with arc angle
    $\Theta\approx 2\arccos\bigl((\rho^2-h^2-\Delta r^2)/(2\rho\Delta r)\bigr)$.
  \item For globally-supported eigenmodes, multiple such shells exist,
    and their projections superpose; constructive interference at
    $r_{\mathrm{proj}} \approx \pi\Reff/(\ell+1)$ generates enhanced
    ring-like or arc-like density contrast.
  \end{enumerate}
\end{theorem}

\begin{proof}
  (i) The intersection of a geodesic sphere $\mathcal{S}_\rho$ with a
  transverse 2-plane in $\Sthree$ is the standard spherical conic
  section. In the large-$\Reff$ limit ($\rho\ll\Reff$), this reduces
  to the Euclidean spherical intersection theorem: the intersection
  is a circle of radius $r_{\mathrm{proj}}=\sqrt{\rho^2-h^2}$ when
  $h<\rho$.

  (ii) When the centre lies within one shell thickness $\Delta r$ of
  $\Sigma_{\mathrm{obs}}$, only a partial arc intersects the slice.
  The arc angle follows from elementary spherical trigonometry.

  (iii) The density-contrast amplitude at $r_{\mathrm{proj}}$ is the
  integral of the harmonic $Y_\ell^{mm'}$ over the shell, which is
  maximised when the projection radius matches the harmonic's preferred
  scale $\pi\Reff/(\ell+1)$.
\end{proof}

\begin{corollary}[Ring $\leftrightarrow$ Arc transition]
  \label{cor:ringarc}
  Whether an $\ell$-mode appears as a ring or as an arc on the observed
  slice depends on the line-of-sight position of the shell centre
  relative to the slice. This explains the simultaneous observation
  of the Giant \emph{Arc} and Big/Giant \emph{Rings} in the same
  redshift shell as projections of different shells with different
  line-of-sight offsets.
\end{corollary}

%% =================================================================
\section{Exact Projected-Chord Law and Asymptotics}
\label{sec:fibonacci}
%% =================================================================

\subsection{Exact Projected-Chord Law}

For the phenomenological size comparison, the more natural observable is
not the geodesic radius itself but the visible chord obtained by
projecting a shell-like compact mode onto a narrow redshift slice. If
$d$ denotes the corresponding projected diameter and $\chi$ the central
geodesic angle on $S^3$, then elementary geometry gives
\begin{equation}
  d = 2\Reff \sin\frac{\chi}{2}.
  \label{eq:chord-basic}
\end{equation}

\begin{proposition}[Exact projected-chord law]
  \label{prop:chordlaw}
  Assume that the characteristic visible diameter associated with a mode
  label $\ell$ is the chord subtending the central angle
  \begin{equation}
    \chi_\ell = \frac{2\pi}{\ell+1}.
  \end{equation}
  Then
  \begin{equation}
    d_\ell = 2\Reff \sin\!\left(\frac{\pi}{\ell+1}\right).
    \label{eq:exactlaw}
  \end{equation}
\end{proposition}

\begin{proof}
  Substitute $\chi=2\pi/(\ell+1)$ into Eq.~\eqref{eq:chord-basic}.
\end{proof}

\begin{remark}[Fibonacci law as an asymptotic corollary]
  For large $\ell$,
  \begin{equation}
    d_\ell
    = 2\Reff \sin\!\left(\frac{\pi}{\ell+1}\right)
    = \frac{2\pi\Reff}{\ell+1} + O\!\left((\ell+1)^{-3}\right).
  \end{equation}
  Hence the earlier v5 scaling is recovered asymptotically if the
  visible subsequence satisfies $\ell_n+1\sim F_n$ or another
  Fibonacci-like coarse-grained hierarchy.
\end{remark}

\subsection{Spectral Separation: Non-Decay vs Scale Quantization}

Two logically distinct ingredients should be separated:
\begin{itemize}[noitemsep]
\item \emph{Compactness $\Rightarrow$ non-decaying global
  correlations.} This follows from Lemma~\ref{lem:notail} and
  Theorem~\ref{thm:nodecay} and does not depend on any Fibonacci or
  mode-selection ansatz.
\item \emph{Projected chord quantization $\Rightarrow$ preferred visible
  diameters.} This follows from Proposition~\ref{prop:chordlaw}. Any
  additional selection of visible mode labels is model-dependent.
\end{itemize}

\section{Phenomenological Identification}
\label{sec:identification}
%% =================================================================

We now offer a phenomenological identification of the observed
structures with explicit mode labels. This section is explicitly labelled
\emph{phenomenological} because it involves empirical calibration of
$\Reff$ and a model-dependent mode assignment, not a first-principles
prediction.

\subsection{Observed Proper Sizes}

All linear sizes used in the comparison are treated as proper transverse
sizes at the mean redshift $z=0.802$ of the Mg\,\textsc{ii}-selected
structures, in the redshift slice $z=0.722$--$0.862$. We use a flat
\LCDM\ background for the distance conversion,
\begin{equation}
  H_0=67\,\mathrm{km\,s^{-1}\,Mpc^{-1}},\quad
  \Omega_m=0.315,\quad \Omega_\Lambda=0.685 .
\end{equation}
For an angular extent $\theta$, the proper transverse size is
\begin{equation}
  d_{\rm proper}=\theta\,D_M(z)/(1+z),
\end{equation}
where $D_M(z)$ is the transverse comoving distance. Varying $H_0$ within
the Planck 2018 one-sigma range changes the common distance scale by less
than one percent, which is smaller than the few-percent residuals in the
phenomenological fit below.

The observed scales used in the fit are
\begin{itemize}[noitemsep]
\item Giant Arc (GA): extent $L_{\mathrm{GA}}=940\,\mathrm{Mpc}$;
\item Giant Ring (GR): diameter $d_{\mathrm{GR}}=840\,\mathrm{Mpc}$;
\item Big Ring (BR): diameter $d_{\mathrm{BR}}=400\,\mathrm{Mpc}$;
\item 2D PSA clustering scale: $\lambda_c=320\,\mathrm{Mpc}$.
\end{itemize}
These values are used as a consistent proper-size set for the present
analysis and should not be interpreted as new measurements of the
Lopez--Clowes structures.

\subsection{Empirical Mode Assignment and One-Parameter Fit}

We replace the earlier Fibonacci-level fit by the exact projected-chord
law of Proposition~\ref{prop:chordlaw}. The empirical assignment is
\begin{equation}
  (\ell_{\mathrm{GA}},\ell_{\mathrm{GR}},\ell_{\mathrm{BR}},\ell_{\mathrm{PSA}})
  =(1,2,6,8).
  \label{eq:ellassign}
\end{equation}
This assignment is inferred from the observed nested morphology and fit
quality; it is not derived from a first-principles selection rule.

Define
\begin{equation}
  s_\ell := 2\sin\!\left(\frac{\pi}{\ell+1}\right),
  \qquad
  d_\ell^{\mathrm{th}} = s_\ell \Reff .
\end{equation}
For observed sizes $d_i^{\mathrm{obs}}$, the ordinary least-squares
estimator for the single free parameter is
\begin{equation}
  \widehat{R}_{\mathrm{eff}}
  = \frac{\sum_i d_i^{\mathrm{obs}} s_i}{\sum_i s_i^2}.
  \label{eq:bestfitR}
\end{equation}
Using $(d_{\mathrm{GA}},d_{\mathrm{GR}},d_{\mathrm{BR}},d_{\mathrm{PSA}})
=(940,840,400,320)\,\mathrm{Mpc}$ gives
\begin{equation}
  \widehat{R}_{\mathrm{eff}} = 474.511\,\mathrm{Mpc}.
\end{equation}

\begin{table}[t]
\centering
\scriptsize
\caption{Single-parameter fit of the exact projected-chord law.}
\label{tab:fitv7}
\begin{tabular}{lcccc}
\toprule
Structure & $\ell$ & $d^{\rm th}$ & $d^{\rm obs}$ & rel. err. \\
\midrule
GA & 1 & 949.022 & 940 & 0.960\% \\
GR & 2 & 821.877 & 840 & 2.157\% \\
BR & 6 & 411.765 & 400 & 2.941\% \\
PSA $\lambda_c$ & 8 & 324.585 & 320 & 1.433\% \\
\bottomrule
\end{tabular}
\end{table}

\begin{remark}[Parameter-free ratios]
The exact law also predicts ratios independent of $\Reff$. The three
nearest checks are
\begin{itemize}[leftmargin=*,noitemsep]
  \item GA/GR: theory $1.1547$; observed $1.1190$.
  \item GR/BR: theory $1.9960$; observed $2.10$.
  \item BR/PSA: theory $1.2686$; observed $1.25$.
\end{itemize}
They provide a parameter-free consistency check at the few-percent level.
\end{remark}

\begin{remark}[Selection-rule caveat]
The fit quality does not by itself derive the visibility of the mode
labels $(1,2,6,8)$. A first-principles selection rule remains open and
must be treated as an explicit unresolved problem rather than as an
established consequence of the model.
\end{remark}

\subsection{The Giant Arc as an $\ell=1$ Mode}

The $\ell=F_1=F_2=1$ assignment deserves comment. On $\Sthree$ the
$\ell=1$ harmonic is a dipolar mode with a single great-circle node:
projected onto a narrow redshift shell, it produces a single
arc-like extended overdensity, not a closed ring. This is
structurally consistent with the observed Giant Arc being an
\emph{arc} rather than a ring~\cite{Lopez2022}. The GR and BR at
$\ell=2,6$ have multiple nodes and can generically project to closed
ring-like structures.



\subsection{Bayesian Model-Comparison Extension}
\label{subsec:bayes-extension}

The one-parameter least-squares fit of Sec.~\ref{sec:identification}
provides a phenomenological consistency check, not a completed model
comparison. A natural next step is to promote the projected-chord law to
a Bayesian model class $M_{S^3}$ and compare it with flat-space, periodic
box, BAO-like, CCC-like, and string-like alternatives.

Let
\begin{equation}
  \mathbf d_{\rm obs}=(L_{\rm GA},d_{\rm GR},d_{\rm BR},\lambda_c)
\end{equation}
be the vector of observed proper sizes. For a model $M$ with parameters
$\theta$, the Gaussian likelihood is
\begin{equation}
\mathcal L(\mathbf d_{\rm obs}|\theta,M)=
\frac{1}{\sqrt{(2\pi)^N|\mathbf C|}}
\exp\!\left[-\frac12\Delta\mathbf d^{T}\mathbf C^{-1}\Delta\mathbf d\right],
\label{eq:likelihood}
\end{equation}
where $\Delta\mathbf d=\mathbf d_{\rm obs}-\mathbf d_{\rm th}(\theta)$.
For the compact-$S^3$ projected-chord model,
\begin{equation}
  d_{\ell}^{\rm th}=2\Reff\sin\!\left(\frac{\pi}{\ell+1}\right),
  \label{eq:chord-bayes}
\end{equation}
with parameters including $\Reff$, the visible mode labels, and nuisance
parameters describing projection and redshift-slice effects. If angular
positions and survey masks are supplied, a topology-informed covariance
may be built from the zonal two-point function,
\begin{equation}
  C_{ij}=\sum_{\ell}P_\ell\,\Xi_\ell(\Theta_{ij})\,W_iW_j
  +\sigma_{\rm noise}^2\delta_{ij},
  \label{eq:covariance}
\end{equation}
where $\Theta_{ij}$ is the angular separation of the pair, $W_i$ denotes
window and selection effects, and $\sigma_{\rm noise}$ includes
measurement and cosmic-variance terms.

Reasonable priors are uniform in $\log\Reff$ over a broad range, e.g.
$10$--$1000\,\mathrm{Mpc}$, and log-uniform for positive low-$\ell$
amplitudes $P_\ell$, optionally regularized by a decaying smoothness
prior. The Bayesian evidence is
\begin{equation}
  Z_M=\int \mathcal L(\mathbf d_{\rm obs}|\theta,M)\,
       \pi(\theta|M)\,d\theta,
  \label{eq:evidence}
\end{equation}
and model comparison proceeds through Bayes factors $B_{ij}=Z_i/Z_j$.
The usual Jeffreys-scale convention treats $\ln B\gtrsim5$ as strong or
decisive support, subject to the adopted priors and likelihood
construction~\cite{Trotta2008,Vardanyan2009}.

\subsubsection{Thermodynamic integration}

One robust numerical route to Eq.~\eqref{eq:evidence} is thermodynamic
integration. Introduce an inverse-temperature parameter $\beta\in[0,1]$
and define the power posterior
\begin{equation}
  p_\beta(\theta)\propto
  \mathcal L(\mathbf d_{\rm obs}|\theta,M)^\beta\,\pi(\theta|M).
  \label{eq:powerpost}
\end{equation}
Then
\begin{equation}
  \log Z_M = \int_0^1
  \left\langle \log\mathcal L(\mathbf d_{\rm obs}|\theta,M)\right\rangle_\beta
  d\beta,
  \label{eq:TI}
\end{equation}
where the expectation value is taken under $p_\beta(\theta)$. In
practice one samples a grid of $\beta_k$ values with MCMC or nested
sampling and performs numerical quadrature over
$\langle\log\mathcal L\rangle_{\beta_k}$. This is the same evidence logic
used in modern cosmological model comparison and in nested-sampling
codes such as PolyChord~\cite{Trotta2008,Handley2015}.

We do not quote a Bayes factor in this paper. The published uLSS size
measurements do not yet provide the joint covariance matrix required for
Eq.~\eqref{eq:likelihood}. Any numerical evidence based on assumed
diagonal errors would therefore be an illustration, not an empirical
result. The purpose of Eqs.~\eqref{eq:likelihood}--\eqref{eq:TI} is to
make the future test explicit: once joint uncertainties, angular
positions, window functions, and absorber selection effects are supplied,
the projected-chord framework becomes a closed statistical model that can
be compared directly with alternatives.

Recent COMPACT Collaboration work provides a methodological precedent
for topology-dependent eigenmode, covariance, and detectability analyses
in cosmic topology. Those templates are not identical to the full
simply-connected $S^3$ shell model studied here, but they demonstrate the
same kind of covariance-based programme required for a decisive test
\cite{COMPACT2024PRL,COMPACT2024IIa}.

\subsection{Additional Projected-Diameter Predictions}
\label{subsec:additional-predictions}

If the empirical mode assignment in Eq.~\eqref{eq:ellassign} reflects a
physical visibility rule rather than an accidental fit, additional
coherent features in the same redshift shell should preferentially occur
near the projected diameters predicted by
Eq.~\eqref{eq:exactlaw}. Using $\Reff=474.511\,\mathrm{Mpc}$ gives the
values in Table~\ref{tab:predictions}. These numbers are not new
observational claims; they are falsifiable targets for deeper DESI,
Euclid, or future Mg\,\textsc{ii} absorber analyses of the same field.

\begin{table}[t]
\centering
\scriptsize
\caption{Projected diameters predicted by the exact chord law for
additional visible modes, using $\Reff=474.511\,\mathrm{Mpc}$.}
\label{tab:predictions}
\begin{tabular}{ccc}
\toprule
Mode $\ell$ & $d_\ell$ [Mpc] & Interpretation \\
\midrule
3  & 671.060 & intermediate shell \\
4  & 557.821 & intermediate shell \\
5  & 474.511 & shell-scale feature \\
7  & 363.175 & below BR scale \\
9  & 293.264 & below PSA scale \\
10 & 267.370 & smaller shell \\
\bottomrule
\end{tabular}
\end{table}

The amplitudes of such features cannot be predicted from the present
size-only fit. A future analysis would need to infer a low-$\ell$ power
spectrum $\{P_\ell\}$ from the two-point function
\begin{equation}
  C(\Theta)=\sum_{\ell\ge 0}P_\ell\Xi_\ell(\Theta),
\end{equation}
combined with a survey window function, noise model, and selection
function. Machine-learning topology classifiers and compactified
simulation methods provide useful methodological precedents, but current
published examples mainly concern toroidal/Euclidean topology or
compactified infinite-space simulations rather than direct nonlinear
structure formation on full $S^3$~\cite{Tamosiunas2024,Racz2018}.

%% =================================================================
\section{Comparison with \LCDM\ and Boundary-Condition Diagnosis}
\label{sec:comparison}
%% =================================================================

\subsection{\LCDM\ Tension Under Finite-Correlation-Length Assumption}

\begin{remark}[Heuristic tension with a finite-correlation-length flat-space picture]
  \label{prop:tension}
  In a spatially flat framework with an effective finite correlation
  length $\xi_0\approx 5$--$10\,\mathrm{Mpc}$, constructing coherent
  ring-like structures of size $\sim 1\,\mathrm{Gpc}$ is generically
  nontrivial and typically strongly suppressed. The precise degree of
  suppression depends on the adopted morphology statistic and on the
  underlying field model.
\end{remark}
\paragraph{Heuristic basis.}
  In Gaussian models with power spectrum $P(k)$, coherent structures of
  physical extent $L$ are increasingly disfavoured once
  $L\gg \xi_{\mathrm{corr}}(P)$. For $L\sim 1\,\mathrm{Gpc}$ and an
  effective late-time correlation scale of order a few to ten
  megaparsecs, the relevant tail behaviour is expected to be strongly
  suppressed. The exact numerical penalty is statistic-dependent, which
  is why we avoid attaching a universal probability estimate here.

\begin{remark}
  We deliberately avoid quoting a specific numerical upper bound here,
  as such a bound requires a specific statistical model of the observed
  uLSS morphology. The qualitative point is that any finite-correlation-length
  assumption generates exponential suppression, which
  Lemma~\ref{lem:notail} shows is inconsistent with compact geometry.
\end{remark}

\subsection{Boundary-Condition Diagnosis of FLAMINGO}

\begin{remark}[Boundary-condition interpretation of FLAMINGO]
  \label{prop:flamingo}
  FLAMINGO-10K~\cite{Sawala2025} and similar large-box cosmological
  simulations use periodic boundary conditions on a cubic box of side
  $L_{\mathrm{box}}\sim$Gpc. Such boundary conditions implement
  a flat torus topology $T^3$, not $\Sthree$. The torus has
  $\pi_1(T^3)=\mathbb{Z}^3$ (violating condition M2), and the
  lowest modes on $T^3$ have wavelength $L_{\mathrm{box}}/n$ for
  integer $n$ with no $SO(4)$-symmetric structure.

  The observation that FLAMINGO fields are consistent with random
  expectations under the 2D PSA of~\cite{Lopez2026} can therefore be
  interpreted as a test of the $T^3$ mode basis used in those
  simulations, not necessarily of all compact topologies. In
  particular, it does not directly test the $\Sthree$ mode basis used
  here.
\end{remark}
\paragraph{Justification.}
  Periodic boundary conditions on $[0,L_{\mathrm{box}}]^3$ identify
  opposite faces and produce $T^3\cong\mathbb{R}^3/\mathbb{Z}^3$.
  The Laplacian on $T^3$ has eigenvalues
  $\lambda_{\mathbf{n}} = (2\pi)^2|\mathbf{n}|^2/L_{\mathrm{box}}^2$
  for $\mathbf{n}\in\mathbb{Z}^3$. These eigenfunctions are
  $\exp(2\pi i\mathbf{n}\cdot\mathbf{x}/L_{\mathrm{box}})$, with no
  analogue of globally-supported low-$\ell$ $\Sthree$ harmonics in the
  isotropic $SO(4)$ sense. The Lopez--Clowes 2D PSA on FLAMINGO fields
  tests this mode basis, which excludes the $\Sthree$ modes.

\begin{remark}[What this means]
  Remark~\ref{prop:flamingo} is an interpretive point about boundary
  conditions: within the present framework, the FLAMINGO result is most
  naturally read as a constraint on the $T^3$ class of simulations
  rather than on the $\Sthree$ mode basis considered here. We are
  therefore offering a topology-explicit interpretation, not an
  empirical claim established by the simulation paper itself.
\end{remark}

\begin{table*}[t]
\centering
\small
\caption{Comparison of $\mathbb{R}^3$, $T^3$ (periodic torus, used in simulations), and compact $\Sthree$ as spatial topologies.}
\label{tab:comparison}
\resizebox{\textwidth}{!}{%
\begin{tabular}{lccc}
\toprule
Property & \LCDM/$\mathbb{R}^3$ & Simulation $T^3$ & Compact $\Sthree$ \\
\midrule
Laplacian spectrum & Continuous & Discrete box modes & Discrete $\ell(\ell+2)/\Reff^2$ \\
Isotropy group & Full $SO(3)$ & Broken to box symmetries & $SO(4)\cong SU(2)_L\times SU(2)_R$ \\
Low-$\ell$ global modes & None ($L^2$) & None ($SO(4)$-incompatible) & $Y_\ell^{mm'}$ for all $\ell$ \\
Fundamental group $\pi_1$ & 0 & $\mathbb{Z}^3$ & 0 \\
Correlations at $\sim$Gpc & Exponentially suppressed & Box-dependent & Non-decaying (Theorem~\ref{thm:nodecay}) \\
Ring-like uLSS & Exponentially suppressed & Not generated in the same mode basis & Structural possibility in the compact-mode picture \\
$\vp$-scale hierarchy & No & No & Emerges only asymptotically from the exact law \\
\bottomrule
\end{tabular}%
}
\end{table*}

\subsection{Limited Observational Consequences}

Under the compact-$\Sthree$ projected-shell hypothesis, the most direct
consequences are the following:

\begin{enumerate}[label=(\arabic*),noitemsep]
\item \textbf{Scale-ratio test.} If additional uLSS features are identified at nearby visible mode labels, their size ratios should follow Eq.~\eqref{eq:exactlaw}; only at large mode number is an approximate $\vp$-hierarchy expected.

\item \textbf{Same-field continuation.} Deeper surveys of the same
  redshift shell may reveal further ring- or arc-like structure at
  smaller radii within the same field.

\item \textbf{Boundary-condition discrimination.} Simulations using
  $T^3$ or $\mathbb{R}^3$ boundary conditions are not expected to probe
  the same compact-mode basis and may therefore underproduce analogous
  systems.
\end{enumerate}

\subsection{Application to DESI Data}
\label{subsec:desi}

The Dark Energy Spectroscopic Instrument (DESI) provides a natural
next-generation test bed for the compact-$S^3$ projected-shell
hypothesis. The DESI Early Data Release includes a public Mg\,\textsc{ii}
absorber value-added catalog with $23{,}921$ absorber entries, and the
DESI DR1 Mg\,\textsc{ii} catalog lists $270{,}529$ absorber entries
\cite{DESIMgIIEDR,DESIMgIIDR1}. The DESI Mg\,\textsc{ii} pipeline was
developed and validated on DESI quasar spectra by Napolitano et al., who
also estimated that the completed five-year survey should contain more
than $8\times10^5$ Mg\,\textsc{ii} absorbers\cite{Napolitano2023}.

For the Lopez--Clowes field, or for equivalent redshift slices selected
from DESI, these data would enable four concrete tests:
\begin{enumerate}[label=(\roman*),noitemsep]
\item a search for same-field structures near the predicted diameters in
  Table~\ref{tab:predictions};
\item construction of a data-driven covariance matrix for
  Eq.~\eqref{eq:covariance}, including survey window and absorber
  selection effects;
\item a Bayesian evidence calculation using Eq.~\eqref{eq:TI}, comparing
  the compact-$S^3$ projected-chord model against flat, periodic-box,
  BAO-like, CCC-like, and string-like templates;
\item a three-dimensional check of whether apparent rings remain
  coherent shell-like structures when DESI redshifts are used rather
  than only projected angular morphology.
\end{enumerate}
If additional coherent features cluster near the predicted scales, this
would strengthen the compact-shell interpretation. Conversely, a null
same-field search together with a Bayes factor disfavoring
$M_{S^3}$ would rule out the present empirical mode assignment at high
confidence. DESI therefore turns the present consistency analysis into a
well-defined falsification programme, but not yet into an established
Bayesian detection.



%% =================================================================
\section{Foundational Consistency Requirements}
\label{sec:foundations}
%% =================================================================

This section records the minimal foundational commitments of the
framework. Its purpose is not to enlarge the observational claim, but to
make explicit which topology, causal structure, entropy principle,
nonlocality statement, and Noether conservation law are being assumed.

\subsection{Topology and Embedding}

The spatial section used in the model is the compact, connected,
orientable, boundaryless, simply connected three-manifold
\begin{equation}
  \Sigma \simeq S^3(\Rcosm) = \{X\in\mathbb R^4: |X|=\Rcosm\}.
\end{equation}
The local effective scale \(\Reff\) used in the projected-shell fit is a
spectral coherence scale in a redshift shell, not the global curvature
radius \(\Rcosm\). The topological status of the spaces relevant to the
construction is summarized as follows:
\begin{table}[t]
\centering
\caption{Topological roles of the basic model spaces.}
\label{tab:topology-foundations}
\begin{ruledtabular}
\begin{tabular}{lcccc}
Space & dim. & compact & connected & $\pi_1$ \\
\hline
$S^1$ & 1 & yes & yes & $\mathbb Z$ \\
$S^2$ & 2 & yes & yes & $0$ \\
$S^3\subset\mathbb R^4$ & 3 & yes & yes & $0$ \\
$\mathbb R^3$ & 3 & no & yes & $0$ \\
\end{tabular}
\end{ruledtabular}
\end{table}
Thus \(S^1\) is compact but not simply connected, \(S^2\) is simply
connected but has the wrong spatial dimension, and \(\mathbb R^3\) has
the correct local dimension but is non-compact and lacks a discrete
global Laplacian spectrum. The model therefore uses \(S^3\) as the
minimal compact simply connected spatial topology; \(\mathbb R^3\)
appears only as the local tangent-space or large-radius approximation.

\subsection{Light Cone and the Value of the Speed of Light}

The spacetime manifold is taken to be
\begin{equation}
  M=\mathbb R\times S^3,
\end{equation}
with Lorentzian metric, in the homogeneous limit,
\begin{equation}
  ds^2=-c^2dt^2+a^2(t)\Rcosm^2 d\Omega_3^2 .
\end{equation}
Massless propagation is determined by the null cone \(ds^2=0\). Locally,
where the tangent space is Minkowskian, this gives
\begin{equation}
  \left|\frac{d\mathbf x}{dt}\right| = c .
\end{equation}
Equivalently, the Maxwell action
\begin{equation}
  S_{\rm em}=-\frac14\int_M d^4x\sqrt{-g}\,F_{\mu\nu}F^{\mu\nu}
\end{equation}
has principal symbol \(g^{\mu\nu}k_\mu k_\nu=0\); its characteristics are
therefore the metric null cones. In this sense the theory fixes the
\emph{origin} of the universal light speed: it is the invariant causal
conversion constant of the local Lorentzian tangent spaces and the
common characteristic speed of all minimally coupled massless fields.
The numerical SI value,
\begin{equation}
  c=299792458\;{\rm m\,s^{-1}},
\end{equation}
is not a new dimensionless prediction of topology; it is the exact
conversion factor fixed by the definitions of the metre and the second.
In natural units used above, this statement is simply \(c=1\).

\subsection{Entropy and the Arrow of Time}

The compact topology does not by itself create an arrow of time. The
microscopic dynamics is assumed to be Hamiltonian or unitary; hence the
fine-grained Gibbs/von-Neumann entropy is conserved,
\begin{equation}
  S_{\rm fine}=-k_B{\rm Tr}\,\rho\ln\rho,\qquad
  \frac{dS_{\rm fine}}{dt}=0 .
\end{equation}
The observed thermodynamic arrow arises from a low-entropy boundary
condition together with coarse graining. If \({\cal C}\) denotes a
coarse-graining map, then
\begin{equation}
  S_{\rm cg}=-k_B{\rm Tr}\,[{\cal C}\rho\ln({\cal C}\rho)]
\end{equation}
can increase under the effective irreversible evolution generated by
tracing over inaccessible microscopic correlations. The model therefore
separates two statements: compact \(S^3\) supplies a finite global mode
basis, while the time arrow is a condition on the state and the
coarse-graining, not on topology alone. For cosmological applications
this is the usual Past-Hypothesis statement: the early slice occupies an
exceptionally small macro-volume in phase space, and later macrostates
occupy larger volumes.

\subsection{Bell Nonlocality and No-Signalling}

Bell nonlocality is included at the level of the quantum state, not as a
superluminal signal on \(S^3\). For two spacelike separated subsystems
with Hilbert space \({\cal H}_A\otimes{\cal H}_B\), an entangled singlet
state gives
\begin{equation}
  E(\mathbf a,\mathbf b)=-\mathbf a\cdot\mathbf b .
\end{equation}
For suitable settings this yields the CHSH value
\begin{equation}
  |S_{\rm CHSH}|=2\sqrt2>2,
\end{equation}
violating the local-hidden-variable bound. At the same time the reduced
state at either wing is independent of the measurement basis chosen at
the other wing, so operational signalling remains excluded. In the
present compact-mode language, Bell correlations are global quantum
state correlations on a common boundaryless slice; they are not mediated
by a signal travelling faster than the null cone fixed above.

\subsection{Noether Symmetry and Conservation Law}

The action is assumed to be invariant under diffeomorphisms of
\(M=\mathbb R\times S^3\) and, for electromagnetism, under local
\(U(1)\) gauge transformations. Variation with respect to the metric
therefore gives the covariant conservation law
\begin{equation}
  \nabla_\mu T^{\mu\nu}=0 .
\end{equation}
In covariant phase-space notation, for a vector field \(\xi\),
\begin{equation}
  J_\xi=\Theta(\phi,{\cal L}_\xi\phi)-\xi\cdot L,
  \qquad dJ_\xi\simeq 0,
\end{equation}
so that on shell
\begin{equation}
  J_\xi\simeq dQ_\xi .
\end{equation}
The Hamiltonian generator satisfies
\begin{equation}
  \delta H_\xi=\int_{\partial\Sigma}
  \left(\delta Q_\xi-\xi\cdot\Theta\right).
\end{equation}
Since \(\Sigma\simeq S^3\) is boundaryless, \(\partial\Sigma=\varnothing\),
and pure-gauge diffeomorphism charges vanish globally:
\begin{equation}
  H_\xi[\Sigma]=0\qquad (\xi\;{\rm gauge}) .
\end{equation}
This is the Noether consistency condition for a closed universe. Local
or quasi-local charges can still be defined on finite subregions with an
artificial boundary, but the total gauge generator on the full compact
slice vanishes, as required by Noether's second theorem.


%% =================================================================
\section{Discussion}
\label{sec:discussion}
%% =================================================================

\subsection{Two-Scale Structure: $\Reff$ vs Global Curvature}

The inferred $\Reff\approx 475\,\mathrm{Mpc}$ at $z=0.8$ is \emph{not}
the global curvature radius of the Universe: the Planck constraint
$|\Omega_k|<5\times 10^{-3}$~\cite{Planck2016} sets
$\Rcosm^{\mathrm{global}}>10\,\mathrm{Gpc}$. The two-scale structure
is as follows:

\begin{itemize}[noitemsep]
\item \emph{Global topological radius} $\Rcosm^{\mathrm{global}}$:
  the actual curvature radius of the cosmological 3-sphere, constrained
  by Planck to $>10\,\mathrm{Gpc}$, governing the CMB power at low $\ell$.
\item \emph{Effective shell curvature scale} $\Reff$: the local
  spectral scale operative within a narrow redshift shell, inferred
  from ring matching to be $474.511\,\mathrm{Mpc}$ at $z=0.8$.
\end{itemize}

The $\Reff$ scale may arise from non-linear collapse and mode-selection
dynamics within the redshift shell, i.e. as a dynamically selected
coherence scale. A detailed derivation would require numerical
simulations on $\Sthree$ (see Open Problems below).

\subsection{Relation to Conformal Cyclic Cosmology}

Penrose's Conformal Cyclic Cosmology (CCC)~\cite{Penrose2014} offers an
alternative mechanism for ring-like features: imprints from
past-aeon black-hole collisions. The compact-$\Sthree$ projected-shell framework is
complementary. Both frameworks predict ring-like structures;
they differ in the predicted correlation between ring positions and
CMB anomalies. CCC predicts correlation with Penrose
circles~\cite{Gurzadyan2013}; our framework predicts enhanced power at
compact projected-shell modes. Future data can distinguish the two.

\subsection{Open Problems}

\begin{enumerate}[label=(\arabic*),noitemsep]
\item \textit{Bayesian evidence}: Construct a covariance matrix for
  $(L_{\mathrm{GA}},d_{\mathrm{GR}},d_{\mathrm{BR}},\lambda_c)$ and compute
  the evidence of the compact-$S^3$ projected-chord model relative to
  flat $\mathbb R^3$, periodic $T^3$, BAO-like, CCC-like, and string-like
  templates.
\item \textit{$N$-body simulations on $\Sthree$}: Numerical structure
  formation on $\Sthree$ with the appropriate boundary conditions
  (not periodic torus) has not been performed. StePS-like compactified
  simulations are a useful precedent, but they are not yet a direct
  nonlinear $S^3$ calculation.
\item \textit{Origin of the $\Reff$ scale}: Why does structure
  formation at $z\approx 0.8$ select $\Reff\approx 475\,\mathrm{Mpc}$?
\item \textit{Mode-selection rule}: The observed set $(1,2,6,8)$ must be
  derived from dynamics, a survey-window selection rule, or a prior on
  the low-$\ell$ power spectrum before the model becomes predictive in a
  first-principles sense.
\item \textit{Same-field prediction test}: Deeper surveys should search
  the same redshift shell for coherent features near the additional
  projected diameters listed in Table~\ref{tab:predictions}.
\item \textit{Interaction with CMB}: Do the same compact shell modes
  generate detectable signatures in the CMB covariance or low-$\ell$
  power spectrum at an observable level?
\end{enumerate}

%% =================================================================
\section{Conclusions}
\label{sec:conclusions}
%% =================================================================

This paper develops a compact-$S^3$ consistency analysis for the
Lopez--Clowes ring-like field near $z\approx 0.8$. The robust analytic
core is independent of any specific mode assignment: compactness gives a
discrete global mode basis, forbids a pure exponential-tail intuition
for globally supported correlations, and naturally allows shell-like
structures whose narrow-slice projections appear as rings or arcs.

For the phenomenological size analysis, we use the exact
projected-chord law
\begin{equation}
  d_\ell = 2\Reff \sin\!\left(\frac{\pi}{\ell+1}\right).
\end{equation}
Earlier Fibonacci scaling appears only as an asymptotic large-$\ell$
limit. With the empirical mode assignment $(1,2,6,8)$ and the best-fit
shell scale $\Reff=474.511\,\mathrm{Mpc}$, the model reproduces the
four observed scales at the percent level.

The paper does not claim that compact topology has been observationally
established, nor that the mode assignment has been derived from first
principles. The narrower claim is that a compact-$S^3$ projected-shell
framework now admits an exact geometric size law and provides a sharply
testable alternative to flat-space intuition for this field. The Bayesian
and same-field tests described above turn this consistency analysis into
a falsifiable programme rather than a completed detection claim.

\begin{acknowledgments}
The manuscript is based on analytical derivations and on comparison
with observational results already available in the cited literature.
\end{acknowledgments}

\section*{Author contributions}
A.P. and B.B. jointly developed the conceptual framework, analytical
arguments, manuscript structure, and interpretation. Both authors
reviewed and approved the submitted version.

\section*{Data availability}
No new observational or simulation datasets were generated in this
study. All data discussed are drawn from published literature and public
preprint sources cited in the reference list.


\appendix
\section{Reproducibility Skeleton for the Projected-Chord Fit}
\label{app:fit-skeleton}

The fit in Table~\ref{tab:fitv7} is reproducible from the following
minimal computation:
\begin{equation}
  s_i=2\sin\!\left(\frac{\pi}{\ell_i+1}\right),\qquad
  \widehat R_{\rm eff}=\frac{\sum_i d_i^{\rm obs}s_i}{\sum_i s_i^2},
\end{equation}
with
\begin{equation}
  (\ell_i)=(1,2,6,8),\qquad
  (d_i^{\rm obs})=(940,840,400,320)\,\mathrm{Mpc}.
\end{equation}
This gives $\widehat R_{\rm eff}=474.511\,\mathrm{Mpc}$. A Bayesian
extension replaces the ordinary least-squares objective by
Eq.~\eqref{eq:likelihood}; no Bayes factor is quoted here because
that requires an observational covariance matrix not yet available from
the published uLSS size estimates.

%% =================================================================
\begin{thebibliography}{99}

\bibitem{Lopez2022}
A.~M.~Lopez, R.~G.~Clowes, and G.~M.~Williger,
A Giant Arc on the sky,
Mon.\ Not.\ Roy.\ Astron.\ Soc.\ \textbf{516}, 1557 (2022),
arXiv:2201.06875.

\bibitem{Lopez2024}
A.~M.~Lopez, R.~G.~Clowes, and G.~M.~Williger,
A Big Ring on the sky,
JCAP \textbf{07}, 055 (2024), arXiv:2402.07591.

\bibitem{Lopez2026}
A.~M.~Lopez and R.~G.~Clowes,
A Giant Ring on the sky,
arXiv:2604.17534 (2026).

\bibitem{Yadav2010}
J.~K.~Yadav, J.~S.~Bagla, and N.~Khandai,
Fractal dimension as a measure of the scale of homogeneity,
Mon.\ Not.\ Roy.\ Astron.\ Soc.\ \textbf{405}, 2009 (2010).

\bibitem{Sawala2025}
T.~Sawala et~al.,
The emperor's new arc: gigaparsec patterns abound in a \LCDM\ universe,
Mon.\ Not.\ Roy.\ Astron.\ Soc.\ \textbf{541}, 1 (2025),
arXiv:2502.03515.

\bibitem{Perelman2002}
G.~Perelman,
The Entropy Formula for the Ricci Flow,
arXiv:math/0211159 (2002).

\bibitem{Perelman2003a}
G.~Perelman,
Ricci Flow with Surgery on Three-Manifolds,
arXiv:math/0303109 (2003).

\bibitem{Perelman2003b}
G.~Perelman,
Finite Extinction Time for the Solutions to the Ricci Flow,
arXiv:math/0307245 (2003).

\bibitem{Planck2016}
Planck Collaboration,
Planck 2015 results.\ XVIII.\ Background geometry and topology,
Astron.\ Astrophys.\ \textbf{594}, A18 (2016).

\bibitem{Camporesi1994}
R.~Camporesi and A.~Higuchi,
Spectral functions and zeta functions in hyperbolic spaces,
J.\ Math.\ Phys.\ \textbf{35}, 4217 (1994).

\bibitem{Berger1971}
M.~Berger, P.~Gauduchon, and E.~Mazet,
\textit{Le spectre d'une vari\'et\'e riemannienne},
Springer Lecture Notes Math.\ \textbf{194} (1971).

\bibitem{Zeckendorf1972}
E.~Zeckendorf,
Repr\'esentation des nombres naturels par une somme de nombres de Fibonacci,
Bull.\ Soc.\ R.\ Sci.\ Li\`ege \textbf{41}, 179 (1972).

\bibitem{Penrose2014}
R.~Penrose,
On the gravitization of quantum mechanics 2: Conformal Cyclic Cosmology,
Found.\ Phys.\ \textbf{44}, 873 (2014).

\bibitem{Gurzadyan2013}
V.~G.~Gurzadyan and R.~Penrose,
On CCC-predicted concentric low-variance circles in the CMB sky,
Eur.\ Phys.\ J.\ Plus \textbf{128}, 2 (2013), arXiv:1302.5162.

\bibitem{CMBS4}
CMB-S4 Collaboration, CMB-S4 Science Book, arXiv:1610.02743 (2016).

\bibitem{LiteBIRD}
LiteBIRD Collaboration,
Probing Cosmic Inflation with the LiteBIRD CMB Polarization Survey,
PTEP \textbf{2023}, 042F01 (2023).

\bibitem{Peebles2022}
P.~J.~E.~Peebles,
Anomalies in physical cosmology,
Ann.\ Phys.\ \textbf{447} (2022), arXiv:2208.05018.


\bibitem{Trotta2008}
R.~Trotta,
Bayes in the sky: Bayesian inference and model selection in cosmology,
Contemp. Phys. \textbf{49}, 71 (2008), arXiv:0803.4089.

\bibitem{Vardanyan2009}
M.~Vardanyan, R.~Trotta, and J.~Silk,
How flat can you get? A model comparison perspective on the curvature of the Universe,
Mon. Not. Roy. Astron. Soc. \textbf{397}, 431 (2009), arXiv:0901.3354.



\bibitem{Handley2015}
W.~J.~Handley, M.~P.~Hobson, and A.~N.~Lasenby,
PolyChord: nested sampling for cosmology,
Mon. Not. Roy. Astron. Soc. \textbf{450}, L61 (2015),
arXiv:1502.01856.

\bibitem{Napolitano2023}
L.~Napolitano et~al.,
Detecting and characterizing Mg~II absorption in DESI survey-validation quasar spectra,
Astron. J. \textbf{166}, 99 (2023).

\bibitem{DESIMgIIEDR}
DESI Collaboration,
Mg~II Absorber Catalog, DESI Early Data Release value-added catalog, version 1.0.

\bibitem{DESIMgIIDR1}
DESI Collaboration,
Mg~II Absorber Catalog, DESI DR1 value-added catalog, version 1.0.
\bibitem{COMPACT2024PRL}
Y.~Akrami et~al. (COMPACT Collaboration),
Promise of future searches for cosmic topology,
Phys. Rev. Lett. \textbf{132}, 171501 (2024).

\bibitem{COMPACT2024IIa}
C.~J.~Copi et~al. (COMPACT Collaboration),
Cosmic topology. Part IIa. Eigenmodes, correlation matrices, and detectability of orientable Euclidean manifolds,
JCAP \textbf{03}, 036 (2024).

\bibitem{Racz2018}
G.~Racz, I.~Szapudi, I.~Csabai, and L.~Dobos,
Compactified cosmological simulations of the infinite universe,
Mon. Not. Roy. Astron. Soc. \textbf{477}, 1949 (2018), arXiv:1711.04959.

\bibitem{Tamosiunas2024}
A.~Tamosiunas et~al. (COMPACT Collaboration),
Cosmic topology. Part IVa. Classification of manifolds using machine learning: a case study with small toroidal universes,
JCAP \textbf{09}, 057 (2024), arXiv:2404.01236.


\bibitem{Bell1964}
J.~S. Bell,
On the Einstein Podolsky Rosen paradox,
Physics \textbf{1}, 195 (1964).

\bibitem{Clauser1969}
J.~F. Clauser, M.~A. Horne, A.~Shimony, and R.~A. Holt,
Proposed experiment to test local hidden-variable theories,
Phys. Rev. Lett. \textbf{23}, 880 (1969).

\bibitem{Noether1918}
E.~Noether,
Invariante Variationsprobleme,
Nachr. Ges. Wiss. Goettingen, Math.-Phys. Kl. 1918, 235 (1918).

\bibitem{IyerWald1994}
V.~Iyer and R.~M. Wald,
Some properties of Noether charge and a proposal for dynamical black hole entropy,
Phys. Rev. D \textbf{50}, 846 (1994).

\bibitem{Wald1984}
R.~M. Wald,
\textit{General Relativity}
(University of Chicago Press, Chicago, 1984).

\end{thebibliography}

\end{document}
